\centerline{\bf \Huge Признаки делимости I}

\begin{flushright}
        \scriptsize{БРБРПАТАПИМ}
\end{flushright}

\problem{\textbf{1}} Начальник написал в тетради число 65349*0712 в качестве примера числа, которое делится: 

а) на 9; б) на 3. (На месте звёздочки когда-то была написана цифра, а теперь там пятно от калмыцкого чая)

Помогите Начальнику восстановить пропущенную цифру. Укажите все возможные варианты!

\problem{\textbf{2}} Запишем подряд цифры от 1 до 9, получим число 123456789. Простое оно или составное? Изменится ли ответ в задаче, если каким-то образом поменять порядок цифр в этом числе?

\problem{\textbf{3}} Покажите, что числа 753, 319, 2457, 6853, 12321, 1234321 являются составными используя признаки делимости!

\problem{\textbf{4}} Найдите числа, которые делятся а) на 9; б) на 4 и заключены между 123456 и 123495.

\problem{\textbf{5}} Напишите наименьшее одиннадцатизначное число, которое делится на 11.

\problem{\textbf{6}} Придумайте а) семизначное; б) десятизначное число, которое делится на 2, 3 и 11, но не
делится на 5 и 9.

\problem{\textbf{7}} Замените звездочки в записи числа 72*4* цифрами так, чтобы это число делилось на 45.
Сколько решений имеет эта задача? Выпишите все решения.